As simulações foram realizadas no ambiente Simulink utilizando a planta
identificada na \Eqref{planta-sensor-1-zn} e os controladores apresentados nas
\Eqref{controlador-pd-paralelo,controlador-ff}. Para aproximar as condições
experimentais, foi adicionado ruído branco à saída da planta e um bloco de
saturação para representar a potência máxima do aquecedor.

Os resultados demonstram que a estratégia composta pelo \gls{pd} e \gls{ff} foi
capaz de acompanhar a referência com baixo erro de seguimento no transitório e
em regime permanente.